\section{\texorpdfstring{\(q\)}{q}-gamma at rational special values}
\label{sec:qgamma}

The \(q\)-gamma function is defined, for \(0<q<1\), by
\[
 \Gamma_q(x)=
 (1-q)^{1-x}\frac{(q;q)_\infty}{(q^x;q)_\infty}
 \qquad\text{\cite[Equation 5.18.4]{dlmf-qgamma}}.
\]
For the complex sector used here, set \(q=\ee^{-t}\),
\(q^x=\ee^{-tx}\), and continue the powers from positive \(t\).
For fixed \(0<x<1\), all product factors remain nonzero when
\(\Re t>0\).  The explicit local branches near \(t=0\) suffice
for the inverse results.

\subsection{An exact modular completion for rational arguments}

\begin{theorem}[Rational-argument multiplication product]
\label{thm:qgamma-completion}
Let \(k\ge2\) be an integer, and define
\[
 \mathcal G_k(t)=\prod_{r=1}^{k-1}\Gamma_{\ee^{-t}}(r/k).
\]
Then, for \(\Re t>0\), with the specified branches,
\begin{equation}\label{eq:qgamma-completion}
 \mathcal G_k(t)=C_k(t)\,\mathcal A_k(Q),
 \qquad Q=\ee^{-4\pi^2/t},
 \qquad \mathcal A_k(Q)=\frac{P(Q)^k}{P(Q^k)},
\end{equation}
where
\begin{equation}\label{eq:qgamma-normalizer}
 C_k(t)=
 \frac{(2\pi)^{(k-1)/2}}{\sqrt{k}}
 \left(\frac{1-\ee^{-t}}{t}\right)^{(k-1)/2}
 \exp\left(\frac{k-1/k}{24}t\right).
\end{equation}
The factor \(C_k\) is holomorphic and nonzero near \(t=0\),
and
\(\mathcal A_k(Q)=1-kQ+\bigO(Q^2)\).
\end{theorem}
\begin{proof}
Splitting positive integers into their residue classes modulo \(k\)
gives
\[
 \prod_{r=1}^{k-1}(q^{r/k};q)_\infty
       =\frac{P(q^{1/k})}{P(q)}.
\]
The powers of \(1-q\) add to \((k-1)/2\).  Consequently
\[
 \mathcal G_k(t)=
 (1-\ee^{-t})^{(k-1)/2}\,
 \frac{P(\ee^{-t})^k}{P(\ee^{-t/k})}.
\]
Apply the \(q=1\) modular formula separately to numerator and
denominator.  Their essential exponentials cancel exactly.
The power constants give
\((2\pi/t)^{(k-1)/2}/\sqrt{k}\);
the regular exponential is
\(\exp((k-1/k)t/24)\).
The dual nome of the denominator is \(Q^k\).
This proves \cref{eq:qgamma-completion,eq:qgamma-normalizer}.
The function \((1-\ee^{-t})/t\) has removable value \(1\)
at the origin, giving the stated analytic root.
Finally, \(P(Q)=1-Q+\bigO(Q^2)\), while \(P(Q^k)=1+\bigO(Q^k)\).
\end{proof}

Thus the \emph{normalized observable}
\[
 \widehat{\mathcal G}_k(t)=\mathcal G_k(t)/C_k(t)
\]
has a flat cusp value \(1\).
Setting
\(y=(1-\widehat{\mathcal G}_k)/k\), the analytic nome inverse
starts \(Q=y+\cdots\), and then
\(t=-4\pi^2/\log Q\).
This is an exact completion of the entire perturbative
\(q\to1\) expansion for these multiplication products.

Every coefficient is arithmetic:
\begin{equation}\label{eq:qgamma-lambda}
 [Q^n]\log\mathcal A_k(Q)
 =-k\sigma_{-1}(n)
   +{\bf1}_{k\mid n}\,\sigma_{-1}(n/k).
\end{equation}
Combining this with \cref{eq:euler-recurrence,eq:psi-lagrange}
computes the inverse to arbitrary order by finite rational
arithmetic.
For example,
\[
 \mathcal A_3(Q)=1-3Q+6Q^3-3Q^4-6Q^7+6Q^9+\cdots .
\]

\subsection{The half-argument theta completion and a sparse inverse}

For \(k=2\), \(\mathcal G_2(t)=\Gamma_{\ee^{-t}}(1/2)\), and
\begin{equation}\label{eq:qgamma-half}
 \Gamma_{\ee^{-t}}(1/2)
 =\sqrt{\pi}\left(\frac{1-\ee^{-t}}t\right)^{1/2}
   \ee^{t/16}\,
   \frac{P(Q)^2}{P(Q^2)},\qquad Q=\ee^{-4\pi^2/t}.
\end{equation}
The Jacobi product gives
\begin{equation}\label{eq:theta-completion}
 \frac{P(Q)^2}{P(Q^2)}
 =\vartheta_4(0,Q)
 =1+2\sum_{n\ge1}(-1)^nQ^{n^2}
 =1-2Q+2Q^4-2Q^9+2Q^{16}-\cdots .
\end{equation}
Here \(\vartheta_4(0,Q)\) is the theta function with nome \(Q\);
the product convention agrees with
\cite[Equation 20.5.4]{dlmf-theta}.
This identity concerns a convergent \(Q\)-series, so its
lacunarity is exact.

Let
\[
 y=\frac{1-\widehat{\mathcal G}_2(t)}2
   =Q-Q^4+Q^9-Q^{16}+\cdots .
\]
Lagrange inversion yields
\begin{align}
 Q(y)={}&y+y^4+4y^7-y^9+22y^{10}
            -13y^{12}+140y^{13}+\bigO(y^{14}),\label{eq:theta-inverse}\\
 \ell(y)=\log(Q(y)/y)
       ={}&y^3+\frac72y^6-y^8+\frac{55}{3}y^9
             -12y^{11}+\frac{455}{4}y^{12}
             +\bigO(y^{13}).\label{eq:theta-log-inverse}
\end{align}
Consequently the principal positive inverse is
\begin{equation}\label{eq:qgamma-normalized-inverse}
 t=\frac{4\pi^2}{\log(1/y)-\ell(y)},\qquad q=\ee^{-t},
\end{equation}
and all other small-nome branches are obtained by the coherent
logarithm lifts in \cref{eq:all-flat-sheets} with \(\nu=1\).

The sparse inverse has a simple support explanation.
Writing
\[
 y=QH(Q),\qquad
 H(Q)=1-Q^3+Q^8-Q^{15}+Q^{24}-\cdots ,
\]
the support of \(H-1\) is \(\{n^2-1:n\ge2\}\).
Its additive monoid is exactly
\[
 \langle3,8\rangle=\{3a+8b:a,b\in\N\}.
\]
The elements \(3,8\) occur.  If \(3\nmid n\), then \(n^2-1\)
is a multiple of \(3\); if \(3\mid n\), then
\(n^2-1=8+(n^2-9)\), also in the monoid.
It follows from \cref{eq:ell-lagrange} that every exponent of
\(\ell\) belongs to \(\langle3,8\rangle\).
The absent powers are therefore structural, rather than
accidental numerical zeros.

\subsection{Reverting the original \texorpdfstring{\(q\)}{q}-gamma value}

The normalized flat observable and the raw \(q\)-gamma value
have different inverses.
For the latter set
\begin{align}
 w&=\log\left(\Gamma_{\ee^{-t}}(1/2)/\sqrt{\pi}\right),\notag\\
 G(t)&=\frac12\log\left(\frac{1-\ee^{-t}}t\right)+\frac{t}{16},
 &L(Q)&=\log\vartheta_4(0,Q).
 \label{eq:raw-gamma-log}
\end{align}
\Cref{eq:qgamma-half} is exactly
\[
 w=G(t)+L(\ee^{-4\pi^2/t}).
\]
The analytic background begins
\begin{equation}\label{eq:gamma-background}
 G(t)=-\frac{3}{16}t+\frac1{48}t^2
       -\frac1{5760}t^4+\frac1{362880}t^6+\cdots ,
 \qquad G'(0)=-\frac3{16}.
\end{equation}
Hence it has an ordinary local inverse
\[
 t_0=G^{-1}(w)=-\frac{16}{3}w+\frac{256}{81}w^2
                 +\bigO(w^3).
\]
Also \(L(Q)=-2Q-2Q^2+\cdots\).
Applying \cref{thm:finite-endpoint} with \(S=4\pi^2\) gives
\begin{equation}\label{eq:raw-gamma-inverse}
 t=t_0+\frac{2}{G'(t_0)}Q_0
       +\sum_{n\ge2}d_n(t_0)Q_0^n,\qquad
 Q_0=\ee^{-4\pi^2/t_0}.
\end{equation}
The first correction is asymptotic to \(-32Q_0/3\).
For every \(n\ge1\),
\begin{equation}\label{eq:gamma-leading-poles}
 d_n(t_0)\sim
 \frac{(-1)^n n^{n-1}}{n!}
 (4\pi^2)^{n-1}\left(\frac{32}{3}\right)^n
 t_0^{2-2n}.
\end{equation}
The exact convergent form is
\[
 t=t_0+t_0^2U\left(t_0,\frac{\ee^{-4\pi^2/t_0}}{t_0^2}\right),
\]
with the uniform error bound \cref{eq:endpoint-tail}.

Thus the raw function has a nonzero sectorial derivative
\(-3\sqrt{\pi}/16\) at its limiting value.
Its algebraic inverse is well defined as an asymptotic core, and
the flat correction supplies the additional terms.
The raw function does not become holomorphic across \(t=0\):
the nonconstant flat factor in \cref{eq:qgamma-half}
still obstructs such an extension.
The ordinary derivative and the essential boundary singularity
are compatible.

For a direct inversion of the raw value rather than its logarithm,
write
\(G_{\rm raw}(t)=\sqrt{\pi}\ee^{G(t)}\).
The first displacement is
\[
 \frac{2G_{\rm raw}(t_0)}{G_{\rm raw}'(t_0)}Q_0
   =\frac{2}{G'(t_0)}Q_0,
\]
so both descriptions agree.

The exponential action can also be expressed directly in the output
deficit.  Let
\[
 z=1-\frac{\Gamma_{\ee^{-t}}(1/2)}{\sqrt\pi},\qquad
 1-\ee^{G(t_0(z))}=z.
\]
Ordinary analytic reversion of this background gives
\[
 t_0(z)=\frac{16}{3}z+\frac{472}{81}z^2+\bigO(z^3),\qquad
 \frac1{t_0(z)}=\frac{3}{16z}-\frac{59}{288}+\bigO(z).
\]
Substitution into the first term of \cref{eq:raw-gamma-inverse}
therefore proves the explicit beyond-all-orders correction
\begin{equation}\label{eq:gamma-deficit-action}
 t-t_0(z)\sim
 -\frac{32}{3}\ee^{59\pi^2/72}
 \ee^{-3\pi^2/(4z)}\qquad (z\downarrow0).
\end{equation}
The transformed action \(3\pi^2/4\) and its nontrivial constant
prefactor both depend on retaining the correct analytic core.

The uniformity region is a strict decay sector for \(t_0\);
a path on which \(\Re(4\pi^2/t_0)\) ceases to grow requires
a different description.

\subsection{Which notion of a critical point is involved?}

There are several distinct questions involving \(q\)-gamma:
varying its argument \(x\) at a fixed base, varying \(q\) toward a
root of unity, removing a regular background, or letting both
variables move.
The special values above concern the base limit \(q\to1\)
at fixed rational arguments.
The ordinary minimum of \(x\mapsto\Gamma_q(x)\), its motion
with \(q\), and a joint critical double scaling are different
problems.  The repository already contains results on that
minimum and the ordinary base expansion \cite{repo-qmono}.
Our theorem supplies exact modular sectors for specific rational
multiplication products and a justified reversion of those sectors.
It does not identify every critical point of a general
\(q\)-special function.
